\documentclass[../comorg-part1.tex]{subfiles}

\begin{document}

\section{General Reference}
\subsection{Size of bytes}
% Shared snippet — \input from reference.tex (and, later, cheatsheet.tex).
% Keep this free of \section headings so the including file controls the level.
\begin{enumerate}
    \item 1 byte = 8 bits
    \item 1KB = $2^{10}$ bytes = 1024 bytes
    \item 1MB = $2^{20}$ bytes = 1,048,576 bytes
    \item 1GB = $2^{30}$ bytes = 1,073,741,824 bytes
\end{enumerate}

\subsection{Memory sizes}
% Shared snippet — \input from reference.tex (and, later, cheatsheet.tex).
\begin{itemize}
    \item Size at each address = 8 bits
    \item Size of address (\textit{assumed}) = 16 bits
    \[
    \text{Size of pointer}_{\text{bytes}} = 8 \times \left\lceil \frac{\text{Address Bits}}{8} \right\rceil
    \]
    \item Size of each register in CPU = 32 bits.\newline
    There are 16 registers in CPU (R0-R12; R13=SP; R14=LP; R15=PC; CSPR)
    \item Memory capacity in BYTES = $2^{\text{Number of address pins}}$
    \item Number of BITS in one data transfer = number of data pins
\end{itemize}

\subsection{Size of numbers}
% Shared snippet — \input from reference.tex (and, later, cheatsheet.tex).
\begin{table}[H]
    \centering
    \begin{tabular}{r c c l}
        \hline
        \textbf{Type} & \textbf{Bytes} & \textbf{Bits} & \textbf{\quad Range} \\
        \hline
        signed char & 1 & 8 & $-128 \to +127$ \\
        unsigned char & 1 & 8 & $0 \to +255$ \\
        short int & 2 & 16 & $-32,768 \to +32,767 \quad (\pm 32\text{KB})$ \\
        unsigned short int & 2 & 16 & $0 \to +65,535 \quad\qquad (64\text{KB})$ \\
        unsigned int & 4 & 32 & $0 \to +4,294,967,295 \quad (4\text{GB})$ \\
        int & 4 & 32 & $-2,147,483,648 \to +2,147,483,647 \ (\pm 2\text{GB})$ \\
        long int & 4 & 32 & $-2,147,483,648 \to +2,147,483,647 \ (\pm 2\text{GB})$ \\
        long long int & 8 & 64 & $-(2^{63}) \to (2^{63})-1$ \\
        float & 4 & 32 & $\pm1.2\times 10^{-38} \to \pm 3.4\times 10^{38} \ (\sim7 \text{ decimals})$ \\
        double & 8 & 64 & $\pm2.3 \times 10^{-308} \to \pm1.7\times 10^{308} \ (\sim 15-17 \text{ decimals})$ \\
        long double & 12 & 96 & $\pm3.4\times 10^{-4932} \to \pm1.1\times 10^{4932} \ (\sim 18-19 \text{ decimals})$ \\
        \hline
    \end{tabular}
\end{table}

\noindent Data types should be aligned to address that is a multiple of the data bus.

\noindent Little Endian (LSB at lowest address; assume this) vs Big Endian (MSB at highest address).


\subsection{Size of instructions and its components/operands}
\begin{enumerate}
    \item One instruction: 32 bits / 4 bytes
    \item Condition flags: 4 bits (Bits 28-31)
    \item Opcode: 8 bits (Bits 20-27)
    \item Register: 4 bits (Since there are 16 possible registers in the processor)
    \item Immediate value: 12 bits (4 rotate right bits + 8 bit immediate value.)
\end{enumerate}

\subsection{Execution Cycle}
Fetch-Decode-Execute $\implies$ PC value is 8 bytes ahead of current instruction being executed.
For PC-related addressing, must have offset of 8bytes.
The max offset is $\pm 4096$ bytes or (\#0xFC) , $\because$ offset field only has 12 bits.


% \subsection{Types of instructions}
% \begin{enumerate}
%     \item MOV/MOVS (destination register) (source register) \newline
%     Copies content from source to destination. MOVS updates status flags after operation.
%     \newline
%     MOV does not touch memory, hence no indirect register methods.
%     \begin{itemize}
%         \item MOV R1, R0
%         \item MOVS R0, R0
%         \item MOV PC, R1
%         \item MOV R1, \#0x100
%     \end{itemize}
%     \item LDR (destination register) [source in memory] \newline
%     Moves 4-byte data from memory to dest register. Cannot have register direct/immediate values.
%     \begin{itemize}
%         \item LDR R1, [R0]: Register R0 stores an address, the processor goes
%         to that address and fetches the data to be copied to R1.
%         \item LDR R1, [R0, \#4]: Register indirect with offset, loads value from address R0+4.
%         \item LDR R1, [R0, \#4]! : Offset w/ autoindexing (pre-index).
%         \item LDR R1, [R0, R2]! : Index register w/ autoindexing (pre-index).
%         \item LDR R1, [R0], \#4 : Offset w/ autoindex (post-index).
%         \item LDR R1, [R0], R2: Index register w/ autoindex (post-index).
%     \end{itemize}
%     \item LDRB (destination) [source]
%     \begin{itemize}
%         \item LDRB R2, [R0] : copy only 1 byte at address in R0 into R2 (byte zero-extends to 32 bits).
%     \end{itemize}
%     \item STR (source register) [destination in memory] \newline
%     Moves 4 byte data from source reg to memory. Cannot have register direct/immediate values.
%     \begin{itemize}
%         \item STR R1, [R0]: Data in R1 copied to the address stored in R0.
%         \item STR R1, [R0, R2]: Register indirect with index register 
%     \end{itemize}
%     \item STRB (source reg) [destination]
%     \begin{itemize}
%         \item R2, [R0, \#1]!: Copy 1 byte in R2 to one address in R0+1.
%     \end{itemize}
%     \item ADD/ADDS (destination reg) (source reg) (source reg / immediate)\newline
%     Affects all flags NZVC. V=1 for signed math, when MSB of operands are same, but MSB of result is different.
%     C=1 for unsigned math, when an extra bit generated that cannot fit into result.\newline
%     Addition is carried out on 32-bit numbers. (registers and adders are 32-bits).
%     \item SUB/SUBS/RSB/RSBS
%     \item MVN/MVNS (destination) (source)\newline
%     Inverts source operand bitwise before moving into destination.
%     \begin{itemize}
%         \item MVN R0, \#0 : R0 = \#FFFFFFFF after execution.
%     \end{itemize}
%     \item AND/ANDS
%     \item ORR/ORRS
%     \item EOR/EORS
%     \item LSL/LSLS/LSR/LSRS
%     \item ASR
%     \item ROR
%     \item RRX
%     \item CMP
%     \item CMN
%     \item TST
%     \item TEQ
% \end{enumerate}

\newpage

\begin{landscape}

\subsection{Instruction Set}
% \pagestyle{empty}

\small 
\renewcommand{\arraystretch}{1.4} 

\begin{longtable}{|p{4cm}|l|p{16cm}|}
\hline
\multicolumn{3}{|l|}{\textbf{\large VisUAL ARM Instruction Set}} \\
\hline
\rowcolor{headercyan} \textbf{Summary} & \textbf{Opcode} & \textbf{Syntax} \\
\hline
\endhead

Move & MOV & \ttfamily MOV\{S\}\{cond\} dest, op1 \{, SHIFT op \#expression\} \newline S=1: NZC, C if shift occurs. \\ \hline
Move Negated & MVN & \ttfamily MVN\{S\}\{cond\} dest, op1 \{, SHIFT op \#expression\} \newline Inverts op1 bitwise (1's comp) before moving. \newline S=1: NZC, C if shift occurs. \\ \hline
LDR Psuedo-Instruction & LDR & \ttfamily LDR\{S\}\{cond\} dest, =expression \\ \hline
Add & ADD & \ttfamily ADD\{S\}\{cond\} dest, op1, op2 \{, SHIFT\_op \#expression\} \newline S=1, updates NZCV. \\ \hline
Add with Carry & ADC & \ttfamily ADC\{S\}\{cond\} dest, op1, op2 \{, SHIFT\_op \#expression\} \newline dest = op1 + op2 + C \newline S=1, updates NZCV. \\ \hline
Subtract & SUB & \ttfamily SUB\{S\}\{cond\} dest, op1, op2 \{, SHIFT\_op \#expression\} \newline S=1, updates NZCV. \\ \hline
Subtract with Carry & SBC & \ttfamily SBC\{S\}\{cond\} dest, op1, op2 \{, SHIFT\_op \#expression\} \newline dest = op1 - op2 - NOT(C) \newline S=1, updates NZCV. \\ \hline
Reverse Subtract & RSB & \ttfamily RSB\{S\}\{cond\} dest, op1, op2 \{, SHIFT\_op \#expression\} \newline dest = op1 - op2 \newline S=1, updates NZCV. \\ \hline
Reverse Subtract with Carry & RSC & \ttfamily RSC\{S\}\{cond\} dest, op1, op2 \{, SHIFT\_op \#expression\} \newline dest = op2 - op1 - NOT(C) \newline S=1, updates NZCV. \\ \hline
Bitwise And & AND & \ttfamily AND\{S\}\{cond\} dest, op1, op2 \{, SHIFT\_op \#expression\} \newline S=1: NZC, C if shift occurs. \\ \hline
Bitwise Exclusive Or & EOR & \ttfamily EOR\{S\}\{cond\} dest, op1, op2 \{, SHIFT\_op \#expression\} \newline S=1: NZC, C if shift occurs. \\ \hline
Bitwise Clear & BIC & \ttfamily BIC\{S\}\{cond\} dest, op1, op2 \{, SHIFT\_op \#expression\} \newline S=1: NZC, C if shift occurs. \\ \hline
Bitwise Or & ORR & \ttfamily ORR\{S\}\{cond\} dest, op1, op2 \{, SHIFT\_op \#expression\} \newline S=1: NZC, C if shift occurs. \\ \hline
Logical Shift Left & LSL & \ttfamily LSL\{S\}\{cond\} dest, op1, op2 \newline S=1: NZC, C gets the last bit shifted out. \\ \hline
Logical Shift Right & LSR & \ttfamily LSR\{S\}\{cond\} dest, op1, op2 \newline S=1: NZC, C gets the last bit shifted out. \\ \hline
Arithmetic Shift Right & ASR & \ttfamily ASR\{S\}\{cond\} dest, op1, op2 \newline MSB unchanged. If S=1, C=LSB. \newline S=1: NZC, C gets the last bit shifted out. \\ \hline
Rotate Right & ROR & \ttfamily ROR\{S\}\{cond\} dest, op1, op2 \newline LSB into MSB. If S=1, C=LSB. \newline S=1: NZC, C gets the last bit shifted out. \\ \hline
Rotate Right and Extend & RRX & \ttfamily RRX\{S\}\{cond\} op1, op2 \newline C into MSB. If S=1, C=LSB. \newline S=1: NZC, C gets the last bit shifted out. \\ \hline
Compare & CMP & \ttfamily CMP\{cond\} op1, op2 \{, SHIFT\_op \#expression\} \newline Test (op1 - op2), same as SUBS without modification \newline Always updates NZCV. \\ \hline
Compare Negated & CMN & \ttfamily CMN\{cond\} op1, op2 \{, SHIFT\_op \#expression\} \newline Tests (op1 + op2), same as ADDS w/o modification \newline Always updates NZCV. \\ \hline
Test Bit(s) Set & TST & \ttfamily TST\{cond\} op1, op2 \{, SHIFT\_op \#expression\} \newline Bitwise AND b/w operands \newline Always updates NZC. \\ \hline
Test Equals & TEQ & \ttfamily TEQ\{cond\} op1, op2 \{, SHIFT\_op \#expression\} \newline Bitwise XOR b/w operands \newline Always updates NZC. \\ \hline
Load Register & LDR & \ttfamily LDR\{B\}\{cond\} dest, [source \{, OFFSET\}] Offset addressing \newline \ttfamily LDR\{B\}\{cond\} dest, [source, OFFSET]! Pre-indexed addressing \newline \ttfamily LDR\{B\}\{cond\} dest, [source], OFFSET Post-indexed addressing \\ \hline
Store Register & STR & \ttfamily STR\{B\}\{cond\} source, [dest \{, OFFSET\}] Offset addressing \newline \ttfamily STR\{B\}\{cond\} source, [dest, OFFSET]! Pre-indexed addressing \newline \ttfamily STR\{B\}\{cond\} source, [dest], OFFSET Post-indexed addressing \\ \hline
Load Multiple Registers & LDM[dir] & \ttfamily LDM[dir]\{cond\} source, \{list of registers\} \\ \hline
Store Multiple Registers & STM[dir] & \ttfamily STM[dir]\{cond\} dest, \{list of registers\} \\ \hline
Branch & B & \ttfamily B\{cond\} target \\ \hline
Branch with Link & BL & \ttfamily BL\{cond\} target \\ \hline
Stop Emulation & END & \ttfamily END\{cond\} \\ \hline
\end{longtable}
\vspace*{\fill}

% \subsection{Memory Access and Clock Cycles per Instruction}
% \begin{table}[H]
% \centering
% \begin{tabular}{|l|c|c|c|}
% \hline
% \textbf{Instruction Type} & \textbf{Memory Accesses (Data + Instr)} & \textbf{Total Accesses} & \textbf{Clock Cycles} \\ \hline

% \textbf{Data Processing} & & & \\
% MOV (Move Register/Immediate) & 1 (Instr) + 0 (Data) & 1 & 1 \\
% MOVEQ (Conditional Move) & 1 (Instr) + 0 (Data) & 1 & 1 (True or False) \\
% MOV LR, PC (Save PC to LR) & 1 (Instr) + 0 (Data) & 1 & 1 \\

% \textbf{Control Flow (Branches modify PC)} & & & \\
% B (Unconditional Branch) & 1 (Instr) + 0 (Data) & 1 & 3 (Pipeline Flush) \\
% BEQ / BCC (Conditional Branch) & 1 (Instr) + 0 (Data) & 1 & 1 (False) or 3 (True) \\
% BL (Branch with Link) & 1 (Instr) + 0 (Data) & 1 & 3 (Pipeline Flush) \\
% BX LR (Branch and Exchange) & 1 (Instr) + 0 (Data) & 1 & 3 (Pipeline Flush) \\
% MOV PC, LR (Return from Branch) & 1 (Instr) + 0 (Data) & 1 & 3 (Pipeline Flush) \\

% \textbf{Single Data Transfer} & & & \\
% LDR (Single Load) & 1 (Instr) + 1 (Data) & 2 & 2 \\
% STR (Single Store) & 1 (Instr) + 1 (Data) & 2 & 2 \\
% LDREQ / STREQ (Conditional) & 1 (Instr) + 0 or 1 (Data) & 1 or 2 & 1 (False) or 2 (True) \\

% \textbf{Block Data Transfer (Stack)} & & & \\
% LDMFD (Load Multiple, Full Desc.) & 1 (Instr) + N (Data) & $N+1$ & $N+2$ \\
% STMFD (Store Multiple, Full Desc.) & 1 (Instr) + N (Data) & $N+1$ & $N+2$ \\
% \hline

% \end{tabular}
% \caption{ARM Memory Accesses vs. Clock Cycles (N = Register List Size)}
% \end{table}

\newpage

\subsection{Memory Access and Clock Cycles per Instruction}
\begin{table}[H]
\centering
\begin{tabular}{|l|c|c|c|}
\hline
\textbf{Instruction Type} & \textbf{Memory Accesses (Data + Instr)} & \textbf{Total Accesses} & \textbf{Clock Cycles} \\ \hline

\textbf{Data Processing (Standard)} & & & \\
MOV, ADD, SUB, AND, ORR, etc. & 1 (Instr) + 0 (Data) & 1 & 1 \\
MOVEQ (Conditional Move) & 1 (Instr) + 0 (Data) & 1 & 1 (True or False) \\
Comparisons (CMP, CMN, TST, TEQ) & 1 (Instr) + 0 (Data) & 1 & 1 \\
Data Proc. with Register Shift (e.g. \texttt{LSL Rs}) & 1 (Instr) + 0 (Data) & 1 & 1 \\

\textbf{Control Flow (Branches modify PC)} & & & \\
B (Unconditional Branch) & 1 (Instr) + 0 (Data) & 1 & 3 (Pipeline Flush) \\
BEQ / BCC (Conditional Branch) & 1 (Instr) + 0 (Data) & 1 & 1 (False) or 3 (True) \\
BL (Branch with Link) & 1 (Instr) + 0 (Data) & 1 & 3 (Pipeline Flush) \\
BX LR (Branch and Exchange) & 1 (Instr) + 0 (Data) & 1 & 3 (Pipeline Flush) \\
MOV PC, LR (Return from Branch) & 1 (Instr) + 0 (Data) & 1 & 3 (Pipeline Flush) \\
Data Proc. writing to PC (e.g. \texttt{ADD PC, ...}) & 1 (Instr) + 0 (Data) & 1 & 3 (Pipeline Flush) \\

\textbf{Single Data Transfer} & & & \\
LDR / STR (Single Load/Store) & 1 (Instr) + 1 (Data) & 2 & 2 \\
LDREQ / STREQ (Conditional) & 1 (Instr) + 0 or 1 (Data) & 1 or 2 & 1 (False) or 2 (True) \\

\textbf{Block Data Transfer (Stack)} & & & \\
LDMFD (Load Multiple, Full Desc.) & 1 (Instr) + N (Data) & $N+1$ & $N+2$ \\
STMFD (Store Multiple, Full Desc.) & 1 (Instr) + N (Data) & $N+1$ & $N+2$ \\
\hline

\end{tabular}
\caption{ARM Memory Accesses vs. Clock Cycles (N = Register List Size)}
\end{table}

\newpage

% \subsection{Memory access per instruction}
% \begin{table}[H]
% \centering
% \begin{tabular}{|l|c|c|}
% \hline
% \textbf{Instruction Type} & \textbf{Memory Accesses (Data + Instruction)} & \textbf{Total Accesses} \\ \hline
% \textbf{Data Processing} & & \\
% MOV (Move Register/Immediate) & 1 (Instr) + 0 (Data) & 1 \\
% \textbf{Control Flow} & & \\
% B / BCC (Branch / Branch Cond.) & 1 (Instr) + 0 (Data) & 1 \\
% \textbf{Single Data Transfer} & & \\
% LDR (Single Load) & 1 (Instr) + 1 (Data) & 2 \\
% STR (Single Store) & 1 (Instr) + 1 (Data) & 2 \\
% \textbf{Block Data Transfer (Stack)} & & \\
% LDMFD (Load Multiple, Full Descending) & 1 (Instr) + N (Data) & $N+1$ \\
% STMFD (Store Multiple, Full Descending) & 1 (Instr) + N (Data) & $N+1$ \\
% \hline
% \end{tabular}
% \caption{ARM Memory Accesses per Instruction (N = Register List Size)}
% \end{table}

\end{landscape}

% \newpage

\newpage

\subsection{Condition Codes}
% Shared snippet — \input from reference.tex (and, later, cheatsheet.tex).
\begin{table}[H]
\centering
\renewcommand{\arraystretch}{1.2}
\begin{tabular}{|l|l|l|l|}
\hline
\textbf{Code} & \textbf{Suffix} & \textbf{Flags} & \textbf{Meaning} \\ \hline
0000 & EQ & $Z = 1$ & Equal \\ \hline
0001 & NE & $Z = 0$ & Not equal \\ \hline
0010 & CS or HS & $C = 1$ & Higher or same, unsigned \\ \hline
0011 & CC or LO & $C = 0$ & Lower, unsigned \\ \hline
0100 & MI & $N = 1$ & Negative \\ \hline
0101 & PL & $N = 0$ & Positive or zero \\ \hline
0110 & VS & $V = 1$ & Overflow \\ \hline
0111 & VC & $V = 0$ & No overflow \\ \hline
1000 & HI & $C = 1$ and $Z = 0$ & Higher, unsigned \\ \hline
1001 & LS & $C = 0$ or $Z = 1$ & Lower or same, unsigned \\ \hline
1010 & GE & $N = V$ & Greater than or equal, signed \\ \hline
1011 & LT & $N \neq V$ & Less than, signed \\ \hline
1100 & GT & $Z = 0$ and $N = V$ & Greater than, signed \\ \hline
1101 & LE & $Z = 1$ and $N \neq V$ & Less than or equal, signed \\ \hline
1110 & AL & Can have any value & Always. This is the default when no suffix is specified. \\ \hline
1111 & NV & & Reserved (unused) \\ \hline
\end{tabular}
\caption*{ARM Bcc Conditions}
\end{table}

\noindent SUBS sets C=1 when no borrow (A - B, where A $>=$ B); clears C=0 if borrow (A - B, where A $<$ B).

\noindent Condition to set the flags:
\begin{itemize}
    \item N: MSB = 1.
    \item Z: All zero.
    \item V: MSB of result $\neq$ MSB of operands, where MSB of operands are the same.
    \item C: Bit carried out = 1.
\end{itemize}


\newpage

\subsection{Instruction Anatomy}
\subsubsection{Data processing instructions}
\begin{figure}[H]
    \centering
    \includegraphics[width=\textwidth]{images/generated/cheatsheet-bf1.svg}
    \caption*{Data Processing with Immediate Value\\ \texttt{XXXCCS Rd, Rs1, \#immediate (rotated)}}
\end{figure}
\begin{figure}[H]
    \centering
    \includegraphics[width=\textwidth]{images/generated/cheatsheet-bf2.svg}
    \caption*{Data Processing with Immediate Shift\\ \texttt{XXXCCS Rd, Rs1, Rs2, \{shft \#\}}}
\end{figure}

\vspace{1em}

\begin{figure}[H]
    \centering
    \includegraphics[width=\textwidth]{images/generated/cheatsheet-bf3.svg}
    \caption*{Data Processing with Register Shift\\ \texttt{XXXCCS Rd, Rs1, Rs2, \{shft Rshift\}}}
\end{figure}

\subsubsection{Branch instructions}
\begin{figure}[H]
    \centering
    \includegraphics[width=\textwidth]{images/generated/cheatsheet-bf4.svg}
    \caption*{Branch Instruction Format; \texttt{B(L)CC Offset}}
\end{figure}
If L=0, then is \texttt{B}. If L=1, then is \texttt{BL}. Offset is distance between current instruct \& target addr. Effectively acts as 26-bit offset.
Branch range: $\pm 32$ Mbytes.

\subsubsection{Load/store instructions}
\begin{figure}[H]
    \centering
    \includegraphics[width=\textwidth]{images/generated/cheatsheet-bf5.svg}
    \caption*{\centering Load/Store Single Word/Byte (Register Offset)\\ \texttt{LDR/STR\{B\}CC Rs, [Rd, offset pre or post]!}}
\end{figure}
\begin{figure}[H]
    \centering
    \includegraphics[width=\textwidth]{images/generated/cheatsheet-bf6.svg}
    \caption*{\centering Load/Store Single Word/Byte (Immediate Offset)\\ \texttt{LDR/STR\{B\}CC Rs, [Rd, offset pre or post]!}}
\end{figure}
\noindent (The maximum offset for register indirect with immediate value is $\pm 4096= \pm 2^{12}$ bytes. The sign is determined by U bit.)\newline
(1/0): P = Pre/post; U = add/subtract; B = byte/32-bit word; W = writeback/no wb; L = load/store;

\begin{figure}[H]
    \centering
    \includegraphics[width=\textwidth]{images/generated/cheatsheet-bf7.svg}
    \caption*{\centering Load/Store Multiple Instructions\\ \texttt{LDM/STM\{mode\}CC Rs\{!\}, \{register list\}}}
\end{figure}
B becomes "don't care" bit; register list bit 0 is R0, bit 15 is R15, 1=include that register.

\begin{table}[H]
    \centering
    \caption*{Summary of Key Instruction Identifiers (Bits 27-25)}
    \begin{tabular}{c l l}
        \toprule
        \textbf{Bits 27--25} & \textbf{Instruction Category} & \textbf{Meaning} \\
        \midrule
        \texttt{000} & Data Processing & Register/Shift Operand ($I=0$) \\
        \texttt{001} & Data Processing & Immediate Operand ($I=1$) \\
        \texttt{010} & Load/Store Single & Register Offset \\
        \texttt{011} & Load/Store Single & Immediate Offset \\
        \texttt{100} & Load/Store Multiple & Block Data Transfer (\texttt{LDM}/\texttt{STM}) \\
        \texttt{101} & Branch              & Branch/Jump Instruction (\texttt{B}/\texttt{BL}) \\
        \bottomrule
    \end{tabular}
\end{table}

\begin{table}[H]
    \centering
    \caption*{Instruction Type (Opcode) Field (Bits 24-21)}
    \begin{tabular}{ll@{\qquad}ll}
        \toprule
        \textbf{Code} & \textbf{Instruction} & \textbf{Code} & \textbf{Instruction} \\
        \midrule
        0000 & AND & 1000 & TST \\
        0001 & EOR & 1001 & TEQ \\
        0010 & SUB & 1010 & CMP \\
        0011 & RSB & 1011 & CMN \\
        0100 & ADD & 1100 & ORR \\
        0101 & ADC & 1101 & MOV \\
        0110 & SBC & 1110 & BIC \\
        0111 & RSC & 1111 & MVN \\
        \bottomrule
    \end{tabular}
\end{table}

\begin{table}[H]
    \centering
    \caption*{Shift Type Encoding (Bits 6 and 5)}
    \begin{tabular}{ll}
        \toprule
        \textbf{Code} & \textbf{Shift type} \\
        \midrule
        00 & LSL \\
        01 & LSR \\
        10 & ASR \\
        11 & ROR/RRX \\
        \bottomrule
    \end{tabular}
\end{table}

% \begin{table}[H]
%     \centering
%     \caption*{Load/Store Multiple: PU Bits (24-23) to Instruction Mapping}
%     \begin{tabular}{ll}
%         \toprule
%         \textbf{PU Bits} & \textbf{Instruction} \\
%         \midrule
%         00 & LDMFA (Pop) \\
%         01 & LDMFD (Pop) \\
%         10 & STMFD (Push) \\
%         11 & STMFA (Push) \\
%         \bottomrule
%     \end{tabular}
% \end{table}

\begin{table}[H]
\centering
\begin{tabular}{cl}
\hline
\textbf{PU\^{}WL Bits} & \textbf{Instruction} \\
\hline
00011 & LDMFA (Pop) \\
01011 & LDMFD (Pop) \\
10010 & STMFD (Push) \\
11010 & STMFA (Push) \\
00010 & STMED (Push) \\
01010 & STMEA (Push) \\
10011 & LDMEA (Pop) \\
11011 & LDMED (Pop) \\
\hline
\end{tabular}
\end{table}

\subsection{Flow control}
\texttt{IF} statement: reverse the condition to reduce jump instructions.\newline
Condtional execution: add the condtion codes to the instruction (e.g. MOVEQ). If using multiple of these, ensure the CSPR unchanged.\newline

\smallskip
\noindent Compound AND: If first condition fail, branch out immediately.\newline
Compound OR: If first cond pass, branch immediately.\newline

\smallskip
\noindent SWITCH with narrow consecutive values, use jump table.\newline
SWTICH with wide range of values, use forked cascade.\newline

\smallskip
\noindent PRE-test loops (WHILE): Test condition at the start, branch out if fail.\newline
POST-test loops (DO-WHILE): Test at the end.\newline

\smallskip
\noindent FOR loops are pre-test loops.

\newpage

\subsection{ASCII table}
% Shared snippet — \input from reference.tex (and, later, cheatsheet.tex).
% Requires \gh (cellcolor shorthand), colortbl, multirow, array.
\begin{center}
\sffamily
\small
\setlength{\tabcolsep}{3.5pt}
\renewcommand{\arraystretch}{1.4}

\begin{tabular}{|c|c|c|c|c| *{8}{c|}}
\hline
% NOTE: pandoc's LaTeX reader silently DROPS any table containing \multirow,
% and likewise chokes on \cline. The corner cell was \multirow{2}{*}{...} with
% \cline{6-13} underneath; flattened here to a plain \multicolumn + \hline so
% the table survives conversion. Visually the corner is no longer merged.
\multicolumn{5}{|c|}{\bfseries Bits: $b_7b_6b_5 \rightarrow$ / $b_4b_3b_2b_1 \downarrow$} &
\gh \bfseries 000 & \gh \bfseries 001 & \gh \bfseries 010 & \gh \bfseries 011 &
\gh \bfseries 100 & \gh \bfseries 101 & \gh \bfseries 110 & \gh \bfseries 111 \\
\hline

\multicolumn{5}{|c|}{} &
\gh \bfseries 0 & \gh \bfseries 1 & \gh \bfseries 2 & \gh \bfseries 3 &
\gh \bfseries 4 & \gh \bfseries 5 & \gh \bfseries 6 & \gh \bfseries 7 \\
\hline

\gh 0 & \gh 0 & \gh 0 & \gh 0 & \gh \bfseries 0 &
NUL & DLE & SP & 0 & @ & P & \`{} & p \\ \hline

\gh 0 & \gh 0 & \gh 0 & \gh 1 & \gh \bfseries 1 &
SOH & DC1 & ! & 1 & A & Q & a & q \\ \hline

\gh 0 & \gh 0 & \gh 1 & \gh 0 & \gh \bfseries 2 &
STX & DC2 & " & 2 & B & R & b & r \\ \hline

\gh 0 & \gh 0 & \gh 1 & \gh 1 & \gh \bfseries 3 &
ETX & DC3 & \# & 3 & C & S & c & s \\ \hline

\gh 0 & \gh 1 & \gh 0 & \gh 0 & \gh \bfseries 4 &
EOT & DC4 & \$ & 4 & D & T & d & t \\ \hline

\gh 0 & \gh 1 & \gh 0 & \gh 1 & \gh \bfseries 5 &
ENQ & NAK & \% & 5 & E & U & e & u \\ \hline

\gh 0 & \gh 1 & \gh 1 & \gh 0 & \gh \bfseries 6 &
ACK & SYN & \& & 6 & F & V & f & v \\ \hline

\gh 0 & \gh 1 & \gh 1 & \gh 1 & \gh \bfseries 7 &
BEL & ETB & ' & 7 & G & W & g & w \\ \hline

\gh 1 & \gh 0 & \gh 0 & \gh 0 & \gh \bfseries 8 &
BS & CAN & ( & 8 & H & X & h & x \\ \hline

\gh 1 & \gh 0 & \gh 0 & \gh 1 & \gh \bfseries 9 &
HT & EM & ) & 9 & I & Y & i & y \\ \hline

\gh 1 & \gh 0 & \gh 1 & \gh 0 & \gh \bfseries 10 &
LF & SUB & * & : & J & Z & j & z \\ \hline

\gh 1 & \gh 0 & \gh 1 & \gh 1 & \gh \bfseries 11 &
VT & ESC & + & ; & K & [ & k & \{ \\ \hline

\gh 1 & \gh 1 & \gh 0 & \gh 0 & \gh \bfseries 12 &
FF & FS & , & < & L & \textbackslash & l & | \\ \hline

\gh 1 & \gh 1 & \gh 0 & \gh 1 & \gh \bfseries 13 &
CR & GS & - & = & M & ] & m & \} \\ \hline

\gh 1 & \gh 1 & \gh 1 & \gh 0 & \gh \bfseries 14 &
SO & RS & . & > & N & \^{} & n & \textasciitilde \\ \hline

\gh 1 & \gh 1 & \gh 1 & \gh 1 & \gh \bfseries 15 &
SI & US & / & ? & O & \_ & o & DEL \\ \hline

\end{tabular}
\end{center}


% \vspace*{\fill}
\newpage

\begin{landscape}

\subsection{Stacks}
% \thispagestyle{empty} 

\vspace*{1em} 

\begin{center}
    \begin{minipage}{\linewidth}
        \centering
        \small % Slightly smaller font to allow narrower columns
        \renewcommand{\arraystretch}{1.8} 
        \setlength{\tabcolsep}{6pt}   
        
        % p{width} used with centering to shrink the table size
        \begin{tabular}{|l|>{\centering\arraybackslash}p{3.5cm}|>{\centering\arraybackslash}p{4cm}|>{\centering\arraybackslash}p{4cm}|>{\centering\arraybackslash}p{3.5cm}|}
        \hline
        \rowcolor{headergray} 
        \textbf{Type} & \textbf{Characteristics} & \textbf{Push Action (Store)} & \textbf{Pop Action (Load)} & \textbf{Access Top Item} \\ 
        \hline
        
        \textbf{FD} & 
        Full Descending \newline Grows Down ($\downarrow$) \newline SP points to: \textbf{Data} & 
        \textbf{Pre-Decrement} \newline \ttfamily STR Rn, [SP, \#-4]! \newline PU\^{}WL: 1010 & 
        \textbf{Post-Increment} \newline \ttfamily LDR Rn, [SP], \#4 \newline PU\^{}WL: 0111 &
        \textbf{No Offset} \newline \ttfamily LDR Rn, [SP] \\ 
        \hline

        \textbf{ED} & 
        Empty Descending \newline Grows Down ($\downarrow$) \newline SP points to: \textbf{Empty} & 
        \textbf{Post-Decrement} \newline \ttfamily STR Rn, [SP], \#-4 \newline PU\^{}WL: 0010 & 
        \textbf{Pre-Increment} \newline \ttfamily LDR Rn, [SP, \#4]! \newline PU\^{}WL: 1111 &
        \textbf{Plus 4} \newline \ttfamily LDR Rn, [SP, \#4] \\ 
        \hline

        \textbf{FA} & 
        Full Ascending \newline Grows Up ($\uparrow$) \newline SP points to: \textbf{Data} & 
        \textbf{Pre-Increment} \newline \ttfamily STR Rn, [SP, \#4]! \newline PU\^{}WL: 1110 & 
        \textbf{Post-Decrement} \newline \ttfamily LDR Rn, [SP], \#-4 \newline PU\^{}WL: 0011 &
        \textbf{No Offset} \newline \ttfamily LDR Rn, [SP] \\ 
        \hline

        \textbf{EA} & 
        Empty Ascending \newline Grows Up ($\uparrow$) \newline SP points to: \textbf{Empty} & 
        \textbf{Post-Increment} \newline \ttfamily STR Rn, [SP], \#4 \newline PU\^{}WL: 0110 & 
        \textbf{Pre-Decrement} \newline \ttfamily LDR Rn, [SP, \#-4]! \newline PU\^{}WL: 1011 &
        \textbf{Minus 4} \newline \ttfamily LDR Rn, [SP, \#-4] \\ 
        \hline
        \end{tabular}
        
        \vspace{15pt}
        \captionof{table}{ARM Stack Operations Summary (SP = R13)}
    \end{minipage}
\end{center}

% \vspace*{\fill} 

\end{landscape}

\subsection{Modular Programming}
Branch to subroutine (both change LR and PC).
\begin{enumerate}
    \item \texttt{BL}.
    \item \texttt{MOV LR PC; B SUBROUTINE}.
\end{enumerate}
Return to calling program (both only change PC).
\begin{enumerate}
    \item \texttt{BX LR}.
    \item \texttt{MOV PC LR}.
\end{enumerate}
\texttt{STMFD SP! \{R6-R7, LR\}}: R6 at lowest memory, LR at highest memory, i.e. SP pointing at R6.\newline
\texttt{LDMFD SP! \{R6-R7, PC\}}: PC takes over saved LR content.\newline
Immediately after executing \texttt{BL Sub1} located at 0x300, LR = 0x304.\newline
Registers:
\begin{itemize}
    \item R0 – R3: Arguments (Inputs) \& Return Values. Volatile / Caller-saved. (Can be safely overwritten).
    \item R4 – R11: Local Variables. Non-volatile / Callee-saved. (MUST be preserved for transparent subroutines).
    \item R11: Frame pointer (use negative offset)
    \item R12 (IP): Intra-procedure scratch register. Volatile. (Can be safely overwritten).
    \item R13 (SP): Stack Pointer. Points to the top of the system stack. (use positive offset)
    \item R14 (LR): Link Register. Holds the return address (PC + 4).
    \item R15 (PC): Program Counter. Holds the address of the currently executing instruction.
\end{itemize}
When pushing LR or extra registers at the start of a nested subroutine, the SP shifts. To recalculate LDR offsets:\newline
Formula: New Offset = Old Offset + (Number of Extra Registers Pushed × 4 bytes)\newline
Example: If a parameter was at [SP, \#24] and you add LR to your push list (STMFD SP!, {R4-R7, LR}), you pushed 1 extra register. New offset = 24 + (1 × 4) = 28.

\vspace{1em}
\begin{verbatim}
LDR R1, [R0], #4
MOV R2, R0
STMFD SP!, {R1, R2}
BL Sub1
\end{verbatim}
Assuming R0=0x100, V1 stored at 0x100, V2 stored at 0x104.\newline
R1 passed in by value, R2 passed in by reference.\newline

\vspace{1em}

\noindent Each recursive call saves the address of the next instruction using LR.
To find number of recursive calls, find the number of repeated addresses saved, that are close to address in PC.



% \newpage

% \begin{figure}[H]
%     \centering
%     \includegraphics[width=0.95\textwidth]{images/condition-codes.png}
%     \caption*{}
%     \label{fig:ldr-instruction}
% \end{figure}

\end{document}
